\section{A scoped threshold and a lower-period shadow}\label{sec:threshold}

\subsection{The first witnessed low-period time sector}

For a smooth projective complex curve $E$ with a finite-order chronological
automorphism $T$, define, in rational Betti cohomology,
\begin{equation}\label{eq:delta-def}
 \delta(E,T)=
 \dim_\Q H^1(E,\Q)-\dim_\Q H^1(E,\Q)^{\langle T\rangle}.
\end{equation}
Over characteristic zero, invariant cohomology is the cohomology of the
smooth quotient, so $\delta$ is the dimension of the sum of nontrivial time
isotypic sectors.

Theorem~\ref{thm:collapse} gives
\begin{equation}\label{eq:delta6}
 \delta(E_6,\tau)=0.
\end{equation}
For comparison, HCS-C20 defines a particular adopted and dynamically
certified period-seven chiral ordered component $E_7$
\cite[Thms.~1 and~3 and Appendix~A]{hcsc20}.  It inherits HCS-C19's explicit
correction of an apparent printed constant
\cite[Prop.~1, Thm.~2, and Cor.~3]{hcsc19}; no
publisher-issued erratum and no classification of the saturated
period-seven scheme are asserted.  Its exact genera are
\begin{equation}\label{eq:c20-genera}
 g(E_7)=8,
 \qquad
 g(E_7/\langle\tau\rangle)=2.
\end{equation}
Consequently,
\begin{equation}\label{eq:delta7}
 \delta(E_7,\tau)=16-4=12.
\end{equation}

\begin{corollary}[Scoped chronology--cohomology threshold]\label{cor:threshold}
For the explicitly source-identified and repository-certified Hamiltonian
H\'enon chiral ordered components considered with $n\le7$, the smallest
period at which at least one component has nontrivial weight-one time
cohomology is $n=7$.
\end{corollary}

\begin{proof}
Gallas's generic class count has no chiral class below period six and one
chiral doublet at period six~\cite[Eqs.~(2)--(7) and Table~1]{gallas2007}.
That doublet's ordered cover
has $\delta=0$ by \eqref{eq:delta6}.  The explicitly defined period-seven
component has $\delta=12$ by \eqref{eq:delta7}.
\end{proof}

The three inputs are displayed without extrapolation in
Table~\ref{tab:scoped-threshold}.
\begin{table}[htbp]
\centering
\begin{tabular}{clcc}
\toprule
Period & Object in the comparison & Genus & $\delta(E_n,\tau)$ \\
\midrule
$n<6$ & no chiral class in the published generic count & -- & -- \\
$6$ & unique source chiral doublet, ordered cover $E_6$ & $1$ & $0$ \\
$7$ & adopted certified HCS-C20 component $E_7$ & $8$ & $12$ \\
\bottomrule
\end{tabular}
\caption{The explicitly scoped low-period chronology comparison.}
\label{tab:scoped-threshold}
\end{table}
The corollary is a first witnessed period in this explicit low-period sample.
It says nothing about all diagonal or non-diagonal components at lower
periods, all period-seven chiral components, or numerical real bounded
orbits at a selected parameter.

\subsection{The quadratic marker is inherited from period one}

The fixed-point equation for \eqref{eq:ham-recurrence} is
\begin{equation}\label{eq:d1}
 D_1(u)=u^2+2u-A=0.
\end{equation}
Besides the period-six reversible marker already recorded in
Section~\ref{sec:source}, the source program supplies a period-seven chiral
quadratic marker
\begin{equation}\label{eq:c7-marker}
 C_7^{\mathrm{mark}}(s_7)=s_7^2-2s_7-A.
\end{equation}
The input markers come from the paragraph following equation~(15) and
equation~(16) of~\cite[p.~3]{endlergallas2006chiral}; the relations below are
direct consequences.

\begin{proposition}[Period-one marker shadow]\label{prop:marker-shadow}
The two marker curves are affine copies of the fixed-point marker:
\begin{equation}\label{eq:marker-alias}
 D_6^{\mathrm{mark}}(s_6)=4D_1(s_6/2),
 \qquad
 C_7^{\mathrm{mark}}(s_7)=D_1(s_7-2).
\end{equation}
Their common function field is $\Q(A,\sqrt{A+1})$.  Their fiber product over
the $A$-line is the reduced union
\begin{equation}\label{eq:fiber-product}
 (s_6-2s_7+4)(s_6+2s_7)=0.
\end{equation}
Its two components meet at
$(A,s_6,s_7)=(-1,-2,1)$, and its normalization is the disjoint union of the
two graphs
\begin{equation}\label{eq:two-graphs}
 s_6=2s_7-4,
 \qquad
 s_6=-2s_7.
\end{equation}
\end{proposition}

\begin{proof}
Both identities in \eqref{eq:marker-alias} follow by substitution.  Solving
either copy of $D_1$ gives the stated quadratic field.  Eliminating $A$ from
$D_6^{\mathrm{mark}}=C_7^{\mathrm{mark}}=0$ gives
\[
 s_6^2+4s_6-4s_7^2+8s_7
 =(s_6-2s_7+4)(s_6+2s_7).
\]
The two lines meet only at $s_7=1,s_6=-2$, for which $A=-1$.  Normalizing a
reduced union of two smooth lines separates their intersection, yielding
\eqref{eq:two-graphs}.
\end{proof}

The proposition concerns $D_6^{\mathrm{mark}}$, not the chiral ordered
curve $E_6$.  It is an obstruction to reading a shared quadratic marker
field as primitive cross-period arithmetic.  It does not prove that every
possible correspondence between the full covers is absent.

\subsection{A restricted clock-divisibility obstruction}

\begin{theorem}[Dominant clock-equivariant maps]\label{thm:divisibility}
Let $X_m$ and $X_n$ be integral varieties with automorphisms $T_m,T_n$ of
exact orders $m,n$.  Suppose $k\in\Z$ and a dominant rational map
$\phi:X_m\dashrightarrow X_n$ satisfies
\begin{equation}\label{eq:equivariance}
 \phi\circ T_m=T_n^k\circ\phi.
\end{equation}
Then
\begin{equation}\label{eq:divisibility}
 \frac{n}{\gcd(n,k)}\mid m,
 \qquad\text{equivalently}\qquad n\mid km.
\end{equation}
If $\gcd(k,n)=1$, then $n\mid m$.
\end{theorem}

\begin{proof}
Iterating \eqref{eq:equivariance} $m$ times gives
\[
 \phi=\phi\circ T_m^m=T_n^{km}\circ\phi.
\]
The image of $\phi$ is dense, so $T_n^{km}$ is the identity as a rational
automorphism of $X_n$.  Since $T_n$ has exact order $n$, the equality
$T_n^{km}=1$ implies $n\mid km$.
\end{proof}

Here clock-faithful means $\gcd(k,n)=1$.  No dominant single-valued
clock-faithful map can therefore connect distinct periods in $\{5,6,7\}$.
The theorem does not exclude non-dominant maps into
boundary fixed loci, $k=0$ clock-forgetting maps, nonfaithful target clocks,
or multivalued algebraic correspondences.  In particular, it is not a
no-Hecke-correspondence theorem.
